\documentclass[12pt,a4paper]{article}

\usepackage[utf8]{inputenc}
\usepackage[T1]{fontenc}
\usepackage[margin=2.5cm]{geometry}
\usepackage{hyperref}
\usepackage{microtype}
\usepackage{parskip}
\usepackage{lmodern}

\hypersetup{colorlinks=true, urlcolor=blue!60!black}

\pagestyle{empty}

\begin{document}

\begin{flushright}
Kundai Farai Sachikonye \\
Auf der Lichtung 19, 82178 Puchheim, Germany \\
\href{https://github.com/fullscreen-triangle}{github.com/fullscreen-triangle} \\
\texttt{kundai.sachikonye@bitspark.com} \\
(+49) 1626386344 \\[6pt]
18 May 2026
\end{flushright}

\vspace{1em}

Dr.\ Seifert \\
Klinik und Poliklinik für Viszeral-, Thorax- und Gefäßchirurgie \\
Universitätsklinikum Carl Gustav Carus Dresden \\
Ref.\ 1776

\vspace{1.5em}

\textbf{Re: Wissenschaftliche*r Mitarbeiter*in --- Tumor Immunology / Gastric Cancer (Ref.\ 1776)}

\vspace{1em}

Dear Dr.\ Seifert,

I am applying for the research associate position in your group (Ref.\ 1776),
studying immune metabolism of tumor-infiltrating T cells in gastric cancer.
My interest in this position is specific: the computational tools I have built
for immunopeptidology --- neoantigen prediction, MHC binding, and immune escape
analysis --- are directly applicable to the T cell biology questions the group
addresses, and they are tools that experimental tumor immunology groups rarely
have in-house. My biochemistry and pharmaceutical biotechnology training provides
the lab skills foundation; the computational capability for T cell epitope analysis
and immune escape modelling is what I bring that is unusual for this role.

\textbf{Computational tools for T cell immunology in gastric cancer.}
The central problem in understanding T cell responses to gastric tumors is
identifying which neoantigens drive productive anti-tumor T cell recognition,
which tumor-derived variants allow the tumor to escape T cell killing, and
which TCR-pMHC interactions are strong enough to sustain effector function
in the metabolically hostile tumor microenvironment.
I have built and deployed all three of these analyses as live tools:

The \href{https://syndrome-seven.vercel.app/tools/neoantigen}{neoantigen predictor}
identifies which somatic mutations produce peptides that are simultaneously
presented on MHC (spectral DFT cosine similarity $\rho_{\mathrm{MHC}} \geq 0.72$)
and immunogenically distinct from self-peptides
($\Delta\rho_{\mathrm{self}} = \rho(\mathrm{WT,self}) - \rho(\mathrm{mut,self}) > 0.15$).
For gastric cancer, this provides a rapid computational screen of the tumor
mutational burden to rank neoantigens by their likelihood of driving productive
CD8$^+$ T cell responses --- before running in vitro binding assays or
organoid co-culture experiments.

The \href{https://syndrome-seven.vercel.app/tools/escape-predictor}{immune escape predictor}
systematically screens all single amino acid substitutions in a target epitope
for variants that retain receptor binding while reducing antibody or T cell
recognition. The escape score $E(v) = -\max\rho(v, \mathrm{Ab}) + w\cdot\rho(v, \mathrm{rec})$
($w=0.80$) identifies substitutions that simultaneously escape immune pressure
and maintain functional interaction. In the gastric cancer context, this is the
computational equivalent of what the organoid system provides experimentally:
a model for why immune-edited tumor clones survive T cell surveillance.

The \href{https://syndrome-seven.vercel.app/tools/mhc-binding}{MHC binding predictor}
operates at $\mathcal{O}(cL\log L)$ complexity, allowing rapid scanning of
patient-specific HLA alleles against the full neoantigen landscape of a tumor.
No database lookup is required --- binding is determined from spectral
complementarity between peptide and HLA physicochemical channels.
Combined with flow cytometry data from tumor-infiltrating T cell subpopulation
analysis, this provides a closed loop: compute predicted binders → validate
presentation by pMHC tetramer staining → correlate with T cell exhaustion
markers (PD-1, TIM-3, LAG-3) in the patient organoid system.

\textbf{Immune metabolism and T cell exhaustion context.}
The T cell metabolic reprogramming in the tumor microenvironment --- the switch
from oxidative phosphorylation to glycolysis, and the downstream effects on
T cell effector function and exhaustion --- has a natural representation in the
disease state framework I have developed.
The \href{https://github.com/fullscreen-triangle/syndrome}{Syndrome} framework
models cellular state as $\mathbf{D}\in[0,1]^8$ across eight oscillator
dimensions including ATP synthesis coherence, genetic expression state, and
circadian regulatory coherence --- three dimensions directly relevant to
T cell metabolic phenotype.
A trajectory on this manifold represents the progression from effector to
exhausted T cell state under metabolic pressure; the sparse $\ell_1$ LP
identifies the minimum intervention (metabolic rescue, checkpoint blockade,
or combination) required to redirect the trajectory.
This framework provides a quantitative language for connecting the metabolic
measurements the group makes (oxygen consumption rate, glycolytic flux,
mitochondrial membrane potential) to T cell functional outcomes.

\textbf{Cancer biology and biochemistry background.}
My BSc in Biochemistry and Cell Biology (Jacobs University Bremen) and
MSc in Pharmaceutical Biotechnology (HAW Hamburg) together cover the practical
biochemistry and cell biology toolkit: cell culture, protein analysis, assay
development, and immunological techniques from training.
At the University Medical Center Groningen (2013) I performed UPLC-MS/MS
shotgun proteomics for cancer biomarker discovery --- work that required
developing and validating sample preparation protocols for tumour-derived
material. In my doctoral research at TUM I worked on mass spectrometry-based
molecular profiling at scale, including federated processing of clinical
multi-omics datasets.
The \href{https://github.com/fullscreen-triangle/lavoisier}{Lavoisier} pipeline
I subsequently built handles spectral proteomics and metabolomics data at
$>$100\,GB scale (98.9\,\% NIST accuracy), applicable to the metabolomic
profiling of T cell metabolic state that the position's research programme
involves.

\textbf{Contribution to the group.}
Experimental tumor immunology groups generate large and structurally complex
datasets --- flow cytometry, scRNA-seq, TCR repertoire sequencing, proteomics,
and organoid imaging. The computational layer for integrating these modalities,
identifying immunogenic neoantigens, predicting escape variants, and modelling
T cell state trajectories is typically not available in-house in a surgical
oncology research group.
The tools I bring address exactly these bottlenecks, and they are available
immediately, deployed, and validated. I would contribute both at the bench
(cell culture, flow cytometry data acquisition and analysis, histological
staining and quantification) and at the analysis layer where the computational
tools create leverage across all experimental outputs.

I hold a BSc in Biochemistry and Cell Biology, an MSc in Pharmaceutical
Biotechnology, and an MSc in Computational Biology (Universität Tübingen).
I am legally resident in Germany with full work authorisation.
English is native; German is conversational. I am available immediately and
willing to relocate to Dresden.

\vspace{1.5em}
Yours sincerely,

\vspace{2.5em}
\textbf{Kundai Farai Sachikonye}

\end{document}
